\section{The thermal determinant and its contact term}
\label{sec:ex-thermal-contact-cancellation}

We now relate the determinant matrices to the stationary integral and
show that replacing the full penalty Hessian by its constraint pullback
does not alter the leading quadratic coefficient. We retain the actual
stationary branch, predictor and tree gauge from the preceding sections.
The group has nonabelian simple real Lie algebra. The finite-volume
stationary Hessian gap and the predictor branch hypotheses remain in
force. We also retain the exact matching of the incoming and outgoing
quadratic forms with the selected bridge; in particular its constraint
map obeys $QW=I-\kappa^{-1}K_c$ for the selected outgoing kinetic
matrix $K_c$. Stationarity and the gap alone do not identify an
independently chosen outgoing matrix. All identities below concern
germs at the trivial field.

\subsection{The Haar factor in tree coordinates}
Let $y$ denote the free tree-gauge coordinates of a stationary point,
and let $j(y)>0$ be the product, over these coordinates, of the Haar
density ratio $J$ of Section~\ref{sec:haar-chart}; it is analytic and
positive near zero, with $j(0)=1$, and by
\eqref{eq:haar-chart-measure} the probability Haar measure in
exponential coordinates has density a positive constant times $J$.
Define the left-increment coordinate map, componentwise, by
\[
 s_y(z)=\log\bigl(\exp(z)\exp(y)^{-1}\bigr),\qquad s_y(y)=0.
\]
It is a local diffeomorphism. Right invariance of probability Haar
measure and the ordinary change-of-variables formula give
\[
                       j(z)=|\det Ds_y(z)|j(s_y(z))
\]
on the common coordinate neighborhood. For completeness, apply both
change-of-variables formulas to any compact subset of that neighborhood.
The two continuous positive densities have equal integrals over every
such set. If they differed at a point, continuity would give a closed
ball of positive Lebesgue measure on which one density is strictly
larger, contradicting that equality. Evaluating at $z=y$ yields
\begin{equation}\label{eq:ex-haar-left-jacobian}
                         j(y)=|\det Ds_y(y)|.
\end{equation}
Thus the Haar density cannot be discarded when the stationary point
moves away from zero.

Let $H$ be the Hessian of the joint action in left coordinates at a
stationary configuration, let $R$ be its fine gauge-tangent matrix,
and let $P$ insert the free tree increments. Gauge invariance and
stationarity give $HR=0$; $H$ is symmetric. The square matrix
$S=[R\ P]$ has $\det(S)^2=1$ in the physical counting bases.
To verify this last assertion, order the fine vertices by tree
height, assigning to a nonroot vertex its upward tree edge and to
a root its coarse scalar coordinate. On the tree edges the gauge
variation is the signed difference between the vertex and its parent,
because those edge variables equal the identity. At a root it is
minus the root gauge parameter. The resulting gauge block is
triangular, with diagonal entries $\pm1$. The free-increment block
is the identity, and the other upper block is zero. Hence the
squared determinant is one, regardless of the free edge variables.

Let $L$ be the gauge-fixing derivative and $\Gamma=LR$ the positive
scalar gauge matrix. Set $C=P^THP$ and $B=LP$. Direct multiplication
gives
\[
 S^T(H+L^TL)S=
 \begin{pmatrix}\Gamma&B\\0&I\end{pmatrix}^{\!T}
 \begin{pmatrix}I&0\\0&C\end{pmatrix}
 \begin{pmatrix}\Gamma&B\\0&I\end{pmatrix}.
\]
Consequently
\begin{equation}\label{eq:ex-gauge-det-factorization}
                  \det(H+L^TL)=\det(\Gamma)^2\det C.
\end{equation}
The tree-chart Hessian $H_{\rm tree}$ is
$Ds_y(y)^T C Ds_y(y)$: the additional second-derivative term in
the chain rule vanishes by stationarity. Combining this with
\eqref{eq:ex-haar-left-jacobian} proves
\[
 j(y)^2\det M_{\rm act}=\det H_{\rm tree}\det\Gamma^2.
\]
The tree gap, positive Haar density, and positive scalar determinant
show that every determinant inside the logarithms is positive near
the origin. Thus
\begin{equation}\label{eq:ex-thermal-determinant}
 \frac12\log\det H_{\rm tree}-\log j(y)
       =\frac12\log\det M_{\rm act}-\log\det\Gamma.
\end{equation}
Adding the inherited loop gives the thermal term of the stationary
Gaussian integral with exactly its original normalization.

\subsection{Replacing the stationary graph by its linear tangent}
Write $\zeta(x)$ for the actual centered stationary representative.
Section~\ref{sec:ex-selected-stationary-linear-data} gives
$\zeta(0)=0$ and, under the stationary tangent input stated there,
$D\zeta(0)x=(Wx,0)$.
Each scalar logarithmic determinant under consideration is invariant
under simultaneous constant color rotation in the coarse, fine and
scalar coordinates. Indeed these rotations act orthogonally; they
conjugate the Hessian and the scalar gauge matrix. The centered-frame
change is equivariant because group multiplication, the nearest-group
map and the logarithm are equivariant under conjugation.

An invariant linear functional on one copy of the Lie algebra is zero.
Its derivative along $\operatorname{Ad}(\exp ta)b$ annihilates every
bracket $[a,b]$. The derived algebra is a nonzero ideal and therefore
is the whole algebra by nonabelian simplicity. Applying this argument
to each coordinate proves that the first derivative of each invariant
scalar function vanishes at zero. The second-derivative chain rule
therefore gives
\[
 D^2\bigl(f(x,\zeta(x))\bigr)_{x=0}[u,v]
   =D^2f(0)[(u,Wu,0),(v,Wv,0)].
\]
This is why the linear external legs used in the Fourier matrix
calculus compute the thermal Hessian of the actual nonlinear branch.
There is no assertion that the nonlinear branches themselves coincide.

Rooted and centered internal increments are related by a triangular
coordinate change with determinant one. On shell the full Hessians
transform by congruence; the constraint-pullback Hessian does so even
off shell. Hence this change preserves their determinants in the
calculations below. Gauge changes preserve the same determinants by
orthogonal congruence. These observations justify working in rooted
coordinates for the contact calculation and centered coordinates for
the uniform estimates.

\subsection{The finite-volume contact identity}
Let $w_x(v)$ be the logarithmic change of the block predictor when an
internal left increment $v$ is applied at background $x$; $w_x(0)=0$.
Let $b(x)$ be the gradient, with respect to this predictor increment,
of the coarse penalty. Write $A(x)=M_{\rm core}(x)$ in rooted
coordinates and define
\[
 S_x(a)_{ij}=\langle a,D_v^2w_x(0)[e_i,e_j]\rangle,
 \qquad Z(x)=S_x(b(x)).
\]
The Hessian chain rule for the penalty says exactly
\[
                         M_{\rm act}(x)=A(x)+Z(x).
\]
It separates the pullback of the penalty Hessian from the gradient
of the penalty paired with the second derivative of its constraint.
At zero $b(0)=0$, and hence $Z(0)=0$.

Define a coarse vector $Y(x)$ by
\[
 \langle a,Y(x)\rangle=\operatorname{tr}\bigl(A(x)^{-1}S_x(a)\bigr).
\]
It is analytic for this fixed finite volume. Gauge covariance follows
directly from this formula: the internal matrices transform by the
same congruence and the predictor output by the coarse adjoint action,
so cyclicity of trace gives the coarse adjoint action on $Y$.
At zero each component of $Y(0)$ is fixed by constant color rotations
and is therefore zero. This also follows by differentiating the adjoint
action: an invariant vector is central, and a nonabelian simple Lie
algebra has zero center.

Let $K_c$ be the coarse kinetic matrix and put $J=-K_c$. The exact
linear bridge identity is
\[
                   QW=I-\kappa^{-1}K_c.
\]
The quadratic penalty at the identity is
$\frac\kappa2|x-z|^2$. Its gradient in $z$ has derivative
$\kappa(z-x)$. Applying the bridge identity therefore gives
\begin{equation}\label{eq:ex-penalty-gradient-derivative}
                              Db(0)=J.
\end{equation}
Along the stationary graph this remains the same derivative, because
its tangent is $W$. In particular $DZ(0)u=S_0(Ju)$.

The identity $\det(A+Z)=\det A\det(1+A^{-1}Z)$ and $Z(0)=0$
give, through second order,
\[
 \log\det(A+Z)-\log\det A
 =\operatorname{tr}(A^{-1}Z)
             -\tfrac12\operatorname{tr}(A^{-1}ZA^{-1}Z)+O(|x|^3).
\]
Since $\operatorname{tr}(A^{-1}Z)=\langle b,Y\rangle$ and
$b(0)=Y(0)=0$, differentiation gives the exact identity
\begin{align}
 D^2(\log\det M_{\rm act}-\log\det M_{\rm core})(0)[u,v]
 &=\langle Ju,DY(0)v\rangle+\langle Jv,DY(0)u\rangle\notag\\
 &\quad-\operatorname{tr}\bigl(A(0)^{-1}S_0(Ju)
                         A(0)^{-1}S_0(Jv)\bigr).
 \label{eq:ex-contact-hessian}
\end{align}
The first two terms include the derivative of the inverse matrix;
it has not been dropped. The vanishing of $Y(0)$ is what removes
the term involving the second derivative of $b$.

There is also a Ward identity
\begin{equation}\label{eq:ex-contact-ward}
                             DY(0)d_ca=0
\end{equation}
for every coarse scalar field $a$. To prove it on the actual branch,
gauge the trivial stationary configuration by $\exp(ta)$ at coarse
sites, then restore tree gauge. Gauge invariance preserves stationarity;
uniqueness in the tree chart identifies the restored configuration with
the selected stationary branch on this curve. The coarse field has
tangent $-d_ca$. Covariance of $Y$ and $Y(0)=0$ show that $Y$ vanishes
along the curve. Differentiation gives
\eqref{eq:ex-contact-ward}. Replacing the branch by its linear tangent
does not alter $DY(0)$, by the first-derivative chain rule. This proves
the identity including the zero modes of the coarse difference.

\subsection{Passing the identity to the common momentum picture}
Fix $L$ for this paragraph. The rooted frame and its change to centered
coordinates are analytic finite stencils. The inverse of $A_0$ is
holomorphic on a fixed strip. These facts construct a common
holomorphic matrix kernel for
$A(x)^{-1}S_x(a)$ and for its square, with two kinds of external
legs: a coarse background leg and a separate coarse test leg $a$.
Indeed express $S_x(a)$ using the second derivative of the local
predictor, conjugate the centered inverse by the triangular change,
and multiply. Every factor has the same finite-volume Fourier
representation; products and inverse derivatives preserve it.
Only a finite bound for this fixed $L$ is required here.

Take normalized traces per coarse site. Finite Fourier inversion
identifies them with averages of the common kernels over the internal
momentum grid. Holomorphy on a smaller strip bounds the Fourier
coefficients exponentially, so these averages converge to the torus
integral, uniformly for real external momentum. The same convergence
holds when an external momentum varies along grid points approaching
a prescribed real momentum: subtract the limiting integral at that
grid point and use its continuity. Thus both the finite-volume
contact identity and its Ward identity pass to all real momenta
along any divergent family of coarse periods.

Denote the resulting mixed background/test trace kernel by
$M_{ij}(p)$, and the second test/test derivative kernel of the trace of the square
by $B_{ij}(p)$. Explicitly, at zero background its finite-volume
bilinear value on test fields $a,b$ is
\[
       2\operatorname{tr}\bigl(A(0)^{-1}S_0(a)A(0)^{-1}S_0(b)\bigr).
\]
Thus $B$ includes both ordered cross terms.
They are analytic near zero. The Ward identity says
\[
                    \sum_i(d_c(-p)a)_iM_{ij}(p)=0.
\]
Fix a component $i=(\nu,\alpha)$ and take $a$ in color $\alpha$
and $p=t e_\nu$. Only that directional difference survives, giving
$q_\nu(-te_\nu)M_{ij}(te_\nu)=0$. The first factor is nonzero
for all sufficiently small nonzero real $t$. Continuity gives
$M_{ij}(0)=0$. Analyticity then implies
$M_{ij}(p)=O(|p|)$ along every ray, not just that coordinate ray.

For a fixed external component $a$, put
$x_j(p)=(-K_c(p)e_a)_j$. Since
$K_c(p)=d_1(-p)^TH_c(p)d_1(p)$ with $H_c$ analytic,
$x_j(p)=O(|p|^2)$. The infinite-volume version of
\eqref{eq:ex-contact-hessian}, including the factor one half in
the thermal determinant, has diagonal density
\[
 \frac12\left[W(p)+W(-p)-\frac12V(p)\right],\qquad
 W(p)=\sum_jx_j(p)M_{aj}(p),\quad
 V(p)=\sum_{i,j}x_i(-p)x_j(p)B_{ij}(p).
\]
The factor $1/2$ in front of $V$ arises because the second derivative
of the trace of the square contains the two equal cross terms.
Now $W$ has ray order at least three and $V$ has ray order at
least four. Their sum therefore has zero quadratic Taylor
coefficient. Passing from a Hessian density to the normalized
quadratic picture introduces one more factor $1/2$ and preserves
this vanishing.

We have proved that the actual and the core one-loop determinant terms,
\[
 \tfrac12\log\det M_{\rm act}-\log\det\Gamma
 \qquad\text{and}\qquad
 \tfrac12\log\det M_{\rm core}-\log\det\Gamma ,
\]
have the same leading quadratic coefficient. By
\eqref{eq:ex-thermal-determinant}, the thermal term of the stationary
integral is the sum of the first of these and the inherited loop
evaluated on the stationary branch. The contact calculation says
nothing about that inherited summand. It is treated in
Section~\ref{sec:inherited-loop-transport}
(Proposition~\ref{prop:inherited-loop-picture}): for the same incoming path,
the same beds and the same stationary bridge, and under the hypotheses
stated there, the common coarse picture of the inherited loop has
order four at zero, so its quadratic Taylor coefficient vanishes. With
that input, the thermal term and the core determinant term have the
same leading quadratic coefficient. This does not say that the
inherited finite-volume Hessian vanishes; only its leading
low-momentum coefficient does.

Combining the exact cancellation for the determinant terms with the
uniform core-versus-rank bound from
Section~\ref{sec:ex-stationary-determinant-prices} bounds the
remaining coefficient difference between the actual and the rank
determinant terms by Cauchy's inequality on the common strip. The
qualitative rooted-frame constants used in the vanishing argument do
not enter that uniform comparison bound.
